\clase{2}{6 de Marzo}{}

\section{Funciones medibles y variables aleatorias}

\subsection{En Análisis}

Sean \((\Omega, \mathcal{F})\) y \((E, \Sigma)\) dos espacios medibles. Una aplicación \(f: \Omega \to E\) se dirá \textit{medible} de \((\Omega, \mathcal{F})\) a \((E, \Sigma)\) (o \((\mcal{F}, \Sigma)\)-medible o medible a secas) si 
\[ f^{-1}(B) \in \mathcal{F} \text{ para todo } B \in \Sigma. \]

Si \((\Omega, \mathcal{F}, \mu)\) es un espacio de medida y \(f: \Omega \to E\) es \((\mcal{F}, \Sigma)\)-medible, la medida \(\mu_f\) sobre \((E, \Sigma)\) dada por la fórmula
\[
\mu_f(B) := \mu \circ f^{-1}(B) = \mu(\underbrace{f^{-1}(B)}_{\in \mathcal{F}})
\]
se denomina la \textit{medida inducida} o \textit{push-forward} de \(\mu\) por \(f\).

\subsection{En Probabilidad}
Una \textit{variable aleatoria} es una aplicación \(X: \Omega \to \mathbb{R}\) medible de \((\Omega, \mathcal{F})\) a \((\mathbb{R}, \mathcal{B}(\mathbb{R}))\), donde:
\begin{itemize}
    \item \((\Omega, \mathcal{F})\) es algún espacio medible,
    \item \(\beta(\mathbb{R})\) es la clase de los \textbf{borelianos} en \(\mathbb{R}\).
\end{itemize}
Es decir, una variable aleatoria es cualquier aplicación medible a valores en \((\mathbb{R}, \beta(\mathbb{R}))\).

\begin{notation}
	Las variables aleatorias se denotan por \(X\) en lugar de \(f\).
\end{notation}

\begin{remark}
    A veces se reemplaza \((\mathbb{R}, \beta(\mathbb{R}))\) por \((\overline{\mathbb{R}}, \beta(\overline{\mathbb{R}}))\) en la definición.
\end{remark}

En general, si $X \colon \Omega \to E$ es medible, vamos a decir que $X$ es un \textit{elemento aleatorio} sobre $E$ (o variable aleatoria a valores en $E$). A lo largo del curso, trabajaremos con elementos aleatorios en donde el espacio $E$ es un espacio topológico.

\begin{definition}
	Un \textit{espacio topológico} es un par $(E, \tau)$ donde:
	\begin{itemize}
		\item $E$ es un espacio;

		\item $\tau$ es una topología sobre $E$.
	\end{itemize}
	Una \textit{topología} sobre $E$ es una clase $\tau$ de subconjuntos de $E$ que satisface:
	\begin{enumerate}[i)]
		\item $\varnothing$ y $E$ pertenecen a $\tau$;

		\item $(U_{i})_{i \in I} \subset \tau \implies \bigcup_{i \in I} U_{i} \in \tau$;

		\item $U_{1}, \dots, U_{k} \in \tau \implies \bigcap_{i=1}^{k} U_{i} \in \tau$.
	\end{enumerate}
	A los conjuntos de la clase $\tau$ se los llama \textit{abiertos} de $E$.
\end{definition}
\medskip
\par En cualquier espacio topológico tenemos una $\sigma$-álgebra canónica:

\begin{definition}
	Dado un espacio $E$ y una clase $\mscr{C}$ de subconjuntos de $E$, definimos la \textit{$\sigma$-álgebra generada por $\mscr{C}$}, y la notamos por $\sigma(\mscr{C})$, como la menor $\sigma$-álgebra de subconjuntos de $E$ que contiene a $\mscr{C}$, i.e, tal que:
	\begin{enumerate}[i)]
		\item $\mscr{C} \subseteq \sigma(\mscr{C})$;

		\item $\mcal{F} \ \sigma$-álgebra, $\mscr{C} \subseteq \mcal{F} \implies \sigma(\mscr{C}) \subseteq \mcal{F}$.
	\end{enumerate}
\end{definition}

\begin{remark}
	Puede verificarse que vale la fórmula
	\[ \sigma(\mscr{C}) \coloneq \bigcap_{\substack{\mcal{F} \ \sigma \text{-álgebra} \\ \mscr{C} \subseteq \mcal{F}}} \mcal{F}. \]
\end{remark}

\begin{definition}
    Dado un espacio topológico \((E, \tau)\), definimos su \textit{\(\sigma\)-álgebra de Borel}, y la notamos \(\mathcal{B}(E)\), como aquella generada por \(\tau\), i.e.
    \[
    \mathcal{B}(E) := \sigma(\tau).
    \]
\end{definition}
\medskip \par
Luego, un \textit{elemento aleatorio} sobre un espacio topológico \(E\) es una aplicación medible \(X: \Omega \to E\) desde un espacio medible \((\Omega, \mathcal{F})\) al espacio medible \((E, \mathcal{B}(E))\). \par

\section{Distribución de una variable aleatoria}

Por último, si $(\Omega, \mcal{F}, \mbb{P})$ es un espacio de probabilidad y $X \colon \Omega \to E$ es $(\mcal{F}, \beta(E))$-medible, entonces la medida $\mbb{P}_{X}$ inducida sobre $(E, \beta(E))$ recibe el nombre de \textit{distribución de $X$} (sobre $(E, \beta(E))$), i.e,
\[ \mbb{P}_{X}(B) \coloneq \mbb{P}(X^{-1}(B)), \quad B \in \beta(E). \]

A los eventos $X^{-1}(B)$ se los suele notar (en probabilidad) como $\{X \in B\}$, i.e, 
\[ \{X \in B\} \coloneq \{\omega \in \Omega \ : \ X(\omega) \in B\} \in \mcal{F}. \]
(Es decir, se suprime el rol de $\omega$ en la notación.)   

\subsection{Diferencia de enfoques}

\begin{itemize}
    \item \textbf{En Análisis}: interesan propiedades de las funciones medibles \(f\) que dependen de su dominio (continuidad, derivabilidad, convexidad, gráfico, etc.).
    \item \textbf{En Probabilidad}: no interesa el dominio de una variable aleatoria \(X\), sino su distribución en el codominio.
\end{itemize}

\section{Función de distribución}

La distribución de una variable aleatoria \(X\) (a valores en \(\mathbb{R}\)) se suele describir mediante su \textbf{función de distribución (acumulada)} \(F_X: \mathbb{R} \to [0,1]\) dada por la fórmula
\[
F_X(x) := \mathbb{P}(X \le x).
\]

\begin{property}
	La función de distribución \(F_X\) satisface:
	\begin{enumerate}[i)]
		\item \(F_X\) es monótona creciente;

		\item $F_{X}(-\infty) := \lim_{x \to -\infty} F_{X}(x) = 0$ y $F_X(+\infty) := \lim_{x \to +\infty} F_X(x) = 1$;

		\item \(F_X\) es continua a derecha, i.e. $F_{X}(x^{+}) \coloneq \lim_{y \to x^{+}} F_{X}(y) = F_{X}(x)$;

		\item $F_{X}(x^{-}) \coloneq \lim_{y \to x^{-}} F_{X}(y) = \mbb{P}(X < x)$;

		\item $\Delta F_X(x) := F_X(x) - F_X(x^-) = \mathbb{P}(X = x)$.
	\end{enumerate}
\end{property}
\medskip \par
Si una función \(F: \mathbb{R} \to [0,1]\) satisface (i), (ii) y (iii) arriba, diremos que se trata de una \textit{función de distribución} (no necesariamente asociada a alguna variable aleatoria \(X\)).
\medskip \par
El siguiente resultado nos dice por qué es suficiente \(F_X\) para describir \(\mathbb{P}_X\).

\begin{theorem}
    Sea \(F\) una función de distribución. Entonces:
    \begin{enumerate}
        \item Existe una variable aleatoria \(X\) (definida en algún espacio de probabilidad) tal que \(F \equiv F_X\).
        \item La distribución de dicha variable es única, i.e.
        \[
        F_{X_1} = F_{X_2} \implies \mathbb{P}_{X_1} = \mathbb{P}_{X_2} \]
		$(\text{si } \mathbb{P}_{X_1} = \mathbb{P}_{X_2} \text{ lo notamos } X_1 \overset{\mathcal{L}}{=} X_2)$.
    \end{enumerate}
\end{theorem}

\begin{corollary}~
    \begin{enumerate}
        \item Dada una probabilidad \(\mathbb{P}\) sobre \((\mathbb{R}, \beta(\mathbb{R}))\), la función \(F: \mathbb{R} \to [0,1]\) definida por \(F(x) := \mathbb{P}((-\infty, x])\) es una función de distribución.

        \item Dada una función de distribución \(F\), existe una única probabilidad \(\mathbb{P}\) sobre \((\mathbb{R}, \beta(\mathbb{R}))\) tal que \(\mathbb{P}((-\infty, x]) = F(x)\) para todo \(x \in \mathbb{R}\). De hecho, \(\mathbb{P} = \mathbb{P}_X\) donde \(X\) es la variable del teorema.
    \end{enumerate}
\end{corollary}
\medskip \par
En particular, existe una correspondencia biunívoca entre
\begin{align*}
	\mscr{A} &\coloneq \{\mbb{P} \ : \ \mbb{P} \text{ medida de probabilidad en } (\R, \beta(\R))\}; \\
	\mscr{B} &\coloneq \{F \ : \ F \text{ función de distribución}\}; \\
	\mscr{C} &\coloneq \{\mbb{P}_{X} \ : \ X \text{ variable aleatoria (a valores en $\R$)}\};  
\end{align*}
dada por:
\begin{align*}
	\alpha \colon \mscr{A} &\to \mscr{B} \\
	\mbb{P} &\mapsto F(x) \coloneq \mbb{P}((-\infty, x]); \\
	& \\
	\beta \colon \mscr{B} &\to \mscr{C} \\
	F &\mapsto \mbb{P}_{X} \text{ donde } F_{X} = F; \\
	& \\
	\gamma \colon \mscr{C} &\to \mscr{A} \\
	\mbb{P}_{X} &\mapsto \mbb{P}_{X}
.\end{align*}
Como consecuencia del Teorema, son inversas
\[ \begin{cases}
	\alpha \text { y } \gamma \circ \beta \\
	\beta \text{ y } \alpha \circ \gamma \\
	\gamma \text{ y } \beta \circ \alpha
\end{cases} \]

Por otro lado, demostraremos únicamente el teorema; la deducción del corolario a partir de éste es sencilla. Además, lo haremos de dos formas:
\begin{enumerate}
    \item \textbf{Forma canónica}: tomamos \((\Omega, \mathcal{F}) := (\mathbb{R}, \beta(\mathbb{R}))\) y \(X := \operatorname{id}_\Omega\) (\(X(\omega) = \omega\)) y definimos \(\mathbb{P}\) sobre \((\Omega, \mathcal{F})\) de manera tal que en \((\Omega, \mathcal{F}, \mathbb{P})\) la variable \(X := \operatorname{id}_\Omega\) cumpla lo buscado.

	\item \textbf{Forma constructiva}: tomamos \((\Omega, \mathcal{F}, \mathbb{P}) := ((0,1), \beta((0,1)), \text{Leb}_{(0,1)})\) y definimos allí una variable \(X\) que cumpla lo buscado.
\end{enumerate}

\section{Tipos de variables aleatorias}

Trabajaremos frecuentemente con dos tipos de variables aleatorias.

\begin{definition}
	Una variable aleatoria \(X\) se dice \textit{discreta} si existe \(S \subseteq \mathbb{R}\) a lo sumo numerable tal que \(\mathbb{P}_X(S) = 1\). (Es decir, si $X$ toma a lo sumo numerables valores, salvo medida nula.)
\end{definition}

Equivalentemente, \(X\) es discreta $\iff \mbb{P}_{X}(\mcal{A}_{X}) = 1$, donde
\[
\mathcal{A}_{X} = \{x \in \mathbb{R} \ : \ \mathbb{P}(X = x) > 0\}
\]
es el conjunto de \textit{átomos} (siempre a lo sumo numerable) de \(X\).

Si \(X\) es discreta, definimos su \textit{función de probabilidad puntual} como la función \(p_{X} : \mathcal{A}_{X} \to [0,1]\) dada por
\[
p_X(x) := \mathbb{P}_{X}(\{x\}). 
\]

\begin{definition}
    Una variable aleatoria \(X\) se dice \textit{absolutamente continua} si existe una función \(f_{X}: \mathbb{R} \to \mathbb{R}_{\ge 0}\) medible Borel tal que
    \[
    F_X(x) = \int_{-\infty}^{x} f_{X}(t)\,dt \quad \text{para todo } x \in \mathbb{R}.
    \]
    En tal caso, la función \(f_{X}\) se llama la \textit{densidad} de \(X\).
\end{definition}

Puede verificarse que \(X\) es absolutamente continua \(\iff\) \(F_X\) es absolutamente continua, i.e. \(F_X\) es derivable en casi todo punto, \(F_X' \in L^1(\mathbb{R})\) y se verifica que
\[
F_X(x) = \int_{-\infty}^{x} F_{X}'(t)\,dt \quad \text{para todo } x \in \mathbb{R}.
\]
Además, en tal caso vale que \(f_X = F_X'\) en casi todo punto.

\begin{remark}
    Existen variables aleatorias que no son ni discretas ni absolutamente continuas, ni una ``mezcla'' de ambas (ver Práctica).
\end{remark}

